% source: https://github.com/pfradin/luadraw/discussions/172#discussioncomment-15442327 (pfradin, block 1)
\documentclass[border=5pt]{standalone}
\usepackage[3d]{luadraw}
\usepackage[svgnames]{xcolor}
\begin{document}
\begin{luadraw}{name=clip3d}
    local g = graph3d:new{
        window={-5,5,-5,5},
        size={10,10},
        viewdir={10,80},
    }
    require 'luadraw_spherical'
    local O = Origin
    Hiddenlines = true; Hiddenlinestyle = "dashed"; g:Linewidth(8)
    g:Define_sphere({radius=5,show=false})
    local C = parallelep(M(-4,-4,-4),8*vecI,8*vecJ,8*vecK)
    --We determine the extremities of the arcs of circles visible on the sphere.
    local AB, CD,EF, GH = {}, {}, {}, {}
    g:DScircle(facet2plane(getfacet(C,1)), {out=AB})
    local A, B = table.unpack(AB)
    g:DScircle(facet2plane(getfacet(C,3)), {out=CD})
    local C1, D = table.unpack(CD)
    g:DScircle(facet2plane(getfacet(C,2)), {out=EF})
    local E, F = table.unpack(EF)
    g:DScircle(facet2plane(getfacet(C,4)), {out=GH})
    local G, H = table.unpack(GH)
    -- now we can calculate the outline of the sphérical part
    local ctr = {A,-4*vecK,B,3,1,-vecK,"ca",O,C1,5,1,"ca", -4*vecJ,D,3,1,-vecJ,"ca", O,E,5,1,"ca", 4*vecK,F,3,1,vecK,"ca", O,G,5,1,"ca",4*vecJ,H,3,1,vecJ,"ca", O,A,5,1,"ca", M(4,3,0),"m",4*vecI,vecI,"c"}
    g:Dpath3d(ctr,"ball color=Chocolate") -- spherical part
    -- now we draw the visible circular facets
    g:Dcircle3d(4*vecK,3,vecK,"fill=Chocolate!50!black"); 
    g:Dcircle3d(4*vecJ,3,vecJ,"fill=Chocolate!50!black")
    g:Dcircle3d(4*vecI,3,vecI,"left color=Chocolate!70!black, right color=Chocolate!50!black")
    -- and the circles with hidden parts
    g:Dspherical()

    g:Show()
\end{luadraw}
\end{document}
